\documentclass[10pt,a4paper]{amsart}
\usepackage[margin=19mm]{geometry}
\usepackage{amsmath,amssymb,mathtools}
\usepackage[scr=rsfs,cal=euler]{mathalfa}
\usepackage{xfrac}
\usepackage[unicode,hidelinks,bookmarks=false]{hyperref}
\newcommand{\ie}{i.e.\ }
\newcommand{\cf}{cf.\ }
\newcommand{\resp}{respectively\ }
\newcommand{\Subsection}{\subsection}
\providecommand{\nobreakdash}{\nobreak\hbox{-}}
\setlength{\emergencystretch}{2em}
\multlinegap0pt
\usepackage{xcolor}
\usepackage{amssymb}
\usepackage{mathtools}
\newtheorem{prop}{Proposition}
\title{Sensitivity and Hamming graphs}
\author{Sara Asensio, Yuval Filmus, Ignacio Garc\'ia-Marco, Kolja Knauer}
\date{Source: arXiv:2505.08951v2 (CC BY 4.0)}
\begin{document}
\maketitle

\noindent\emph{Source and licence.} Sensitivity and Hamming graphs. Version arXiv:2505.08951v2 (CC BY 4.0), distributed by arXiv under the Creative Commons Attribution 4.0 International licence, http://creativecommons.org/licenses/by/4.0/, as recorded in the arXiv licence metadata for this exact version and shown on the abstract page https://arxiv.org/abs/2505.08951v2. Full licence notice: \texttt{licenses/arxiv-2505.08951-CC-BY-4.0.txt}.\par\smallskip

\noindent\textit{Editorial scope.} Counted unit: the article's prop prop:FirstConstruction (source0.tex lines 109--113). The article's complete original proof of it is delivered with it (source0.tex lines 114--160, one \texttt{proof} environment of 47 lines). No mathematics is written, repaired or re-derived: the statement and the proof are the article's own lines, in the article's own notation, and no retained line is substituted or re-flowed.  Retained uncounted as the source-local context the proof needs: lines 162--166.  One declared adaptation: exactly 1 ancillary strings inside the retained spans are omitted (image-inclusion commands and, where the source repeats a label, the repeated label definitions) and no other line differs from the source; the figure environments, with their placement, captions and labels, are kept, and the package records the omitted strings in fidelity receipts bound to the affected bodies.  No retained line contains a citation, so the wrapper carries no bibliography and no bibliography entry.  The retained lines contain the article's own diagram code, which the wrapper typesets with the standard tikz packages; no image file is delivered and the package declares no asset dependency.  Editorial formatting only, in addition: an AMS \texttt{amsart} preamble that restates the article's own declarations and macros used by the retained lines, namely the environments \texttt{prop} on separate counters, together with the standard packages those lines need.  The retained environments print in this document as Proposition 1 in order of appearance, whereas the article prints the same items with the numbering of its own full document.  The complete native source of the article as distributed in the arXiv e-print for this exact version is bundled unchanged as \texttt{sources/178079/source0.tex}.
 \begin{figure}[!ht]
 \centering
 	\caption{The construction in Proposition \ref{prop:FirstConstruction} with $m=4$ and $n=3$} 
 	\label{fig:counterexample3}
 \end{figure}

\begin{prop}[{\sf Imbalanced partitions of degree $1$}]
\label{prop:FirstConstruction}
For all integers $m\geq 1, n\geq 2$ there exists a partition $\Pi$ of $H(n,m)$ into $m$ sets with maximum degree $\Delta(\Pi)\leq 1$ and imbalance \[\iota(\Pi)= \begin{cases} m-2 & \text{if } m \text{ is even,}\\
m-1 & \text{if } m \text{ is odd.}\end{cases}\]
\end{prop}

\begin{proof}
    For every $i\in\{0,1,\dots,m-1\}$, let
\[
 S_i = \bigcup_{k=0}^{n-1} \left\{ x  \ \#\ b \ \#\ 0^k : |x| + k + 1 = n,\, b \neq 0,\, \Sigma(x) + \left\lfloor \frac{b+1}{2} \right\rfloor \equiv i \pmod{m} \right\}\ ,
\]
where $\#$ represents concatenation and $\Sigma(x)$ is the sum of the entries of $x$. Moreover, when $x$ is empty  we consider that it sums to $0$. Consider the following induced subgraphs of $H(n,m)$:
    \begin{itemize}
        \item $H_0$ is the induced subgraph on $S_0\cup\{0^n\}$.
        \item $H_i$ is the induced subgraph on $S_i$ for all $i\in\{1,2,\dots,m-1\}.$
    \end{itemize}

For an illustration of the case 
 $m=4$ and $n=3$,
see Figure~\ref{fig:counterexample3}.

    
    These $m$ induced subgraphs constitute a partition of $H(n,m)$. 
Let us show first that each $S_i$ has maximum degree at most ~$1$. Suppose that $x\ \#\ b\ \#\ 0^k,\ y\ \#\ c\ \#\ 0^\ell \in S_i$ are neighbors. We consider two cases:
\begin{enumerate}
    \item $k = \ell$. In this case $x = y$ and so $\lfloor \frac{b+1}{2} \rfloor = \lfloor \frac{c+1}{2} \rfloor$. Given $b$ there is at most one other option for $c$.
    \item $k \neq \ell$, without loss of generality $k > \ell$. In this case $y= x \ \#\ b \ \#\ 0^{k-\ell-1}$. So
    \[
     \Sigma(x) + \left\lfloor \frac{b+1}{2} \right\rfloor \equiv \Sigma(x) + b + \left\lfloor \frac{c+1}{2} \right\rfloor \pmod{m} \ \Rightarrow\ 
     \left\lfloor \frac{b+1}{2} \right\rfloor - b \equiv \left\lfloor \frac{c+1}{2} \right\rfloor \pmod{m}\ .
    \]
    The left-hand side is $-\left\lceil \frac{b+1}{2} \right\rceil+1$, and so it ranges from $-\left\lceil \frac{m}{2}\right\rceil +1$ to $0$, which is to say from $\left\lfloor \frac{m}{2} \right\rfloor + 1$ to $m$.
    The right-hand side ranges from $1$ to $\left\lfloor\frac{m}{2} \right\rfloor$.
    So the two sides cannot be equal.
\end{enumerate}
The only vertex of $H(n,m)$ not covered by the sets $S_i$ is $0^n$. Its neighbors are of the form $0^i \ \#\ b \ \#\ 0^{n-i-1}$, where $b \neq 0$. In this case $x = 0^i$ and so $\Sigma(x) + \left\lfloor \frac{b+1}{2} \right\rfloor = \left\lfloor \frac{b+1}{2} \right\rfloor \in \{1, \ldots, \left\lfloor \frac{m}{2} \right\rfloor\}$. Therefore $0^n$ does not have neighbors in $S_0$ and hence $H_0$ maintains the maximum degree at most~$1$.

To complete the analysis of the construction, let us compute the sizes of the sets $S_i$:
\[
 |S_i| = (m-1)\sum_{j=1}^{n-1} m^{j-1} + \left|\left\{ b \neq 0 : \left\lfloor \frac{b+1}{2} \right\rfloor = i\right\}\right|.
\]
The first summand equals $m^{n-1} - 1$. When $i = 1$, the second summand is $2$. Hence the partition is imbalanced. In fact, it is possible to compute the exact imbalance of this construction from the value of the second summand in the previous expression.

On the one hand, when $m$ is even, the second summand takes the following values, which lead to an imbalance of $m-2$:

\begin{itemize}
    \item $0$ when $i=0$,
    \item $2$ for $i\in\{1,2,\dots,\frac{m}{2}-1\}$,
    \item $1$ when $i=\frac{m}{2}$, and
    \item $0$ for $i\in\{\frac{m}{2}+1,\dots,m-1\}$.
\end{itemize}

On the other hand, when $m$ is odd, the imbalance is equal to $m-1$. The difference compared to the previous case is that $|\{b\neq 0 : \lfloor\frac{b+1}{2}\rfloor=i\}|$ equals $2$ when $i\in\{1,2,\dots,\frac{m-1}{2}\}$ and $0$ otherwise.\end{proof}

\end{document}
